\chapter{Limits, Integrals, and Approximate Identities}
\label{ch:limits}

\begin{quote}
\itshape
Everything so far has used the word ``tends'' without saying in which sense. This chapter supplies
the analysis: three inequivalent notions of convergence, one theorem that licenses the interchange
of limit and integral, and a complete proof that any kernel family with unit mass concentrating at a
point samples continuous functions at that point.
\end{quote}

\section{Three modes of convergence}

Let $f_{n},f$ be functions on $\R$ and let $T_{n},T$ be distributions.

\begin{definition}\label{def:modes}
\begin{enumerate}
\item $f_{n}\to f$ \emph{pointwise} if $f_{n}(x)\to f(x)$ for each fixed $x$.
\item $f_{n}\to f$ \emph{uniformly} if $\sup_{x}|f_{n}(x)-f(x)|\to0$.
\item $T_{n}\to T$ \emph{in the distributional sense} if $\ip{T_{n}}{\varphi}\to\ip{T}{\varphi}$ for
each fixed test function $\varphi$.
\end{enumerate}
\end{definition}

The three are not ordered by strength in the way one might hope, and on an unbounded domain none of
them controls the behaviour of integrals.

\begin{example}[uniform convergence does not preserve the integral]\label{ex:uniform}
Let $f_{n}(x)=n^{-1}\mathbf{1}_{[0,n]}(x)$. Then $\sup_{x}|f_{n}(x)|=1/n\to0$, so $f_{n}\to0$
uniformly, while
\begin{equation}
\int_{\R}f_{n}(x)\dd x=\frac{1}{n}\cdot n=1
\qquad\text{for every }n .
\end{equation}
The functions become shorter and longer at the same rate; the area is constant.
\end{example}

\begin{example}[pointwise convergence is compatible with a delta limit]\label{ex:delta-limit}
Let $g_{n}(x)=n\,\mathbf{1}_{(0,1/n)}(x)$. For fixed $x\neq0$ one has $g_{n}(x)=0$ as soon as
$n>1/x$, so $g_{n}\to0$ pointwise. Nevertheless, for $\varphi$ continuous at $0$,
\begin{equation}
\ip{g_{n}}{\varphi}=n\int_{0}^{1/n}\varphi(x)\dd x\longrightarrow\varphi(0),
\end{equation}
that is, $g_{n}\to\delta$ distributionally. The same sequence tends to $0$ pointwise and to
$\delta$ in the distributional sense; the two statements are about different objects.
\end{example}

\begin{remark}
For the Cauchy kernel the pointwise statement has to be qualified once more: $K_{\ep}(x)\to0$ for
$x\neq0$, but $K_{\ep}(0)=1/(\pi\ep)\to\infty$. Statements of the form ``the kernel tends to zero
everywhere'' are therefore false, and the exception at the origin is precisely the carrier of the
mass.
\end{remark}

\section{When may a limit enter an integral?}

\begin{theorem}[dominated convergence]\label{thm:DCT}
Let $f_{n}\to f$ almost everywhere. Suppose there is a single $g\in\Lone$ with
$|f_{n}(x)|\le g(x)$ almost everywhere, for all $n$. Then
\begin{equation}
\int_{\R}f_{n}\longrightarrow\int_{\R}f .
\end{equation}
\end{theorem}

The hypothesis that is usually misread is the phrase \emph{a single} $g$: a different bound for each
$n$ is not enough.

\begin{proposition}[no common dominator for the Cauchy kernels]\label{prop:nodominator}
There is no $g\in\Lone$ with $K_{\ep}\le g$ for all $\ep\in(0,1)$.
\end{proposition}

\begin{proof}
Suppose such a $g$ exists. Since $K_{\ep}(x)\to0$ for each $x\neq0$, dominated convergence gives
$\int K_{\ep}\to0$. But $\int K_{\ep}=1$ for every $\ep$, a contradiction.
\end{proof}

\begin{remark}[what the test function buys]
The family $e^{-\ep|t|}$ is dominated by the constant $1$, which is not integrable; so
Theorem~\ref{thm:DCT} does not apply to $\int e^{-\ep|t|}\dd t$ over the whole line. Multiply by a
test function, however, and
\begin{equation}
\big|e^{-\ep|t|}\varphi(t)\big|\le|\varphi(t)|\in\Lone ,
\end{equation}
a bound independent of $\ep$ and integrable. Dominated convergence then gives
$\ip{e^{-\ep|\cdot|}}{\varphi}\to\ip{1}{\varphi}$ for every $\varphi\in\Sw$.

The introduction of test functions is therefore not a formality: it is the operation that replaces a
non-integrable dominating function by an integrable one. This single observation is the technical
reason the whole theory is built on pairings.
\end{remark}

\section{Approximate identities}

\begin{definition}\label{def:approxid}
Let $k\in\Lone$ with $\int_{\R}k=1$. For $\ep>0$ put
\begin{equation}
k_{\ep}(x)=\frac{1}{\ep}\,k\!\left(\frac{x}{\ep}\right).
\label{eq:dilate-k}
\end{equation}
The family $(k_{\ep})_{\ep>0}$ is called an \emph{approximate identity}.
\end{definition}

\begin{lemma}[mass and tails]\label{lem:mass-tail}
For every $\ep>0$, $\int_{\R}k_{\ep}=1$. For every $a>0$,
\begin{equation}
\int_{|x|>a}k_{\ep}(x)\dd x=\int_{|u|>a/\ep}k(u)\dd u\longrightarrow0
\qquad(\ep\downarrow0).
\label{eq:tail}
\end{equation}
\end{lemma}

\begin{proof}
The first statement is the substitution $x=\ep u$. For the second, the same substitution converts the
left side into the tail of the fixed integrable function $k$ over the expanding set
$\{|u|>a/\ep\}$; the tails of an integrable function vanish.
\end{proof}

\begin{theorem}[sampling]\label{thm:sampling-general}
Let $(k_{\ep})$ be an approximate identity and let $\varphi$ be bounded and continuous at $0$. Then
\begin{equation}
\int_{\R}k_{\ep}(x)\varphi(x)\dd x\longrightarrow\varphi(0).
\label{eq:sampling-general}
\end{equation}
\end{theorem}

\begin{proof}
Because the total mass is $1$, the difference is
\begin{equation}
\int_{\R}k_{\ep}\varphi-\varphi(0)=\int_{\R}k_{\ep}(x)\big[\varphi(x)-\varphi(0)\big]\dd x .
\end{equation}
Let $\eta>0$. Continuity at $0$ gives $a>0$ such that $|x|<a$ implies
$|\varphi(x)-\varphi(0)|<\eta$. Splitting at $|x|=a$ and using $k_{\ep}\ge0$ and
$\|k_{\ep}\|_{1}=1$,
\begin{equation}
\left|\int_{\R}k_{\ep}\varphi-\varphi(0)\right|
\le\eta\int_{|x|<a}k_{\ep}+\int_{|x|\ge a}k_{\ep}\,2\|\varphi\|_{\infty}
\le\eta+2\|\varphi\|_{\infty}\int_{|x|\ge a}k_{\ep}.
\end{equation}
Now fix $a$ and let $\ep\downarrow0$; the last term tends to $0$ by \eqref{eq:tail}. Hence the
superior limit is at most $\eta$, for every $\eta>0$.
\end{proof}

\begin{remark}[the structure of the estimate]
The proof uses two different mechanisms in two different regions: near the origin the smallness of
$\varphi(x)-\varphi(0)$ comes from \emph{continuity}, and far from the origin the smallness of the
remaining mass comes from \emph{integrability}. The parameter $a$ is fixed before $\ep$ is sent to
zero, and $\eta$ is sent to zero last. Reversing this order of quantifiers is the standard way to
produce a wrong proof.
\end{remark}

For the specific Cauchy kernel a change of variable gives a shorter argument, and the comparison is
instructive.

\begin{proposition}\label{prop:cauchy-sub}
For bounded continuous $\varphi$,
\begin{equation}
\int_{\R}K_{\ep}(x)\varphi(x)\dd x=\frac{1}{\pi}\int_{\R}\frac{\varphi(\ep u)}{1+u^{2}}\dd u
\longrightarrow\varphi(0).
\end{equation}
\end{proposition}

\begin{proof}
Substitute $x=\ep u$. The integrand is dominated by
$\|\varphi\|_{\infty}/(\pi(1+u^{2}))$, which does not involve $\ep$ and is integrable; for each fixed
$u$, $\varphi(\ep u)\to\varphi(0)$. Dominated convergence applies, and
$\frac{1}{\pi}\int(1+u^{2})^{-1}\dd u=1$.
\end{proof}

The two proofs differ in what is held fixed. In Theorem~\ref{thm:sampling-general} the kernel varies
and the test function is fixed; in Proposition~\ref{prop:cauchy-sub} the kernel is rescaled into a
fixed weight and the variation is transferred to the test function. Only the second one produces a
dominating function, which is why the first one needs the two-region estimate instead.

\section{Convolution and the identity element}

\begin{definition}
For suitable $k,f$,
\begin{equation}
(k*f)(x)=\int_{\R}k(x-y)f(y)\dd y=\int_{\R}k(y)f(x-y)\dd y .
\end{equation}
\end{definition}

\begin{theorem}\label{thm:conv}
Let $(k_{\ep})$ be an approximate identity and let $f$ be bounded and continuous at $x$. Then
$(k_{\ep}*f)(x)\to f(x)$. If $f$ is bounded and uniformly continuous, the convergence is uniform;
if $f\in L^{p}(\R)$ with $1\le p<\infty$, then $k_{\ep}*f\to f$ in $L^{p}$.
\end{theorem}

\begin{proof}[Proof of the pointwise statement]
Apply Theorem~\ref{thm:sampling-general} to the function $y\mapsto f(x-y)$, which is bounded and
continuous at $y=0$:
\begin{equation}
(k_{\ep}*f)(x)-f(x)=\int_{\R}k_{\ep}(y)\big[f(x-y)-f(x)\big]\dd y\longrightarrow0 . \qedhere
\end{equation}
\end{proof}

\begin{proposition}[$\delta$ is the identity]\label{prop:delta-identity}
Formally and, once $\delta$ is a distribution, rigorously,
\begin{equation}
(\delta*f)(x)=\int_{\R}\delta(y)f(x-y)\dd y=f(x).
\end{equation}
\end{proposition}

\begin{proof}
By the defining property of $\delta$ applied to the test function $y\mapsto f(x-y)$.
\end{proof}

So the phrase ``approximate identity'' is literal: convolution with $k_{\ep}$ approximates the
identity map on functions, and $\delta$ is that map.

\section{How much of the structure is needed}

\begin{proposition}[general sufficient conditions]\label{prop:general-kernels}
Let $(k_{\ep})_{\ep>0}\subset\Lone$ satisfy
\begin{enumerate}
\item $\displaystyle\int_{\R}k_{\ep}=1$ for every $\ep$;
\item $\displaystyle\sup_{\ep}\|k_{\ep}\|_{1}<\infty$;
\item $\displaystyle\int_{|x|>a}|k_{\ep}(x)|\dd x\to0$ for every $a>0$.
\end{enumerate}
Then $\ip{k_{\ep}}{\varphi}\to\varphi(0)$ for every bounded continuous $\varphi$.
\end{proposition}

\begin{proof}
Repeat the proof of Theorem~\ref{thm:sampling-general} with $k_{\ep}$ replaced by $|k_{\ep}|$; the
near term is bounded by $\eta\sup_{\ep}\|k_{\ep}\|_{1}$ and the far term by
$2\|\varphi\|_{\infty}\int_{|x|\ge a}|k_{\ep}|$, which tends to $0$.
\end{proof}

Nonnegativity is thus a convenience, not a necessity.

\begin{proposition}[mass can escape]\label{prop:escaping}
Let $k_{n}=\tfrac12\mathbf{1}_{[n-1,n+1]}$. Then $\int_{\R}k_{n}=1$ for every $n$, but
$\ip{k_{n}}{\varphi}\to0$ for every $\varphi\in\Sw$; that is, $k_{n}\to0$, not $\delta$.
\end{proposition}

\begin{proof}
For $\varphi\in\Sw$ one has $\sup_{|x-n|\le1}|\varphi(x)|=O(n^{-2})$ by rapid decay, hence
\begin{equation}
|\ip{k_{n}}{\varphi}|\le\frac{1}{2}\int_{n-1}^{n+1}|\varphi(x)|\dd x=O(n^{-2})\longrightarrow0 .
\end{equation}
\end{proof}

\begin{remark}[why this does not contradict mass conservation]
The total mass is conserved, but it escapes to infinity. In the topology of $\Sw'$ the only
continuous tests are against rapidly decaying functions, and the constant function $1$ is not one of
them. Consequently ``total mass'' is not a continuous functional on $\Sw'$, and a family may conserve
mass while converging to zero. This is not a pathology to be avoided; it is the reason the theory
must be formulated with test functions rather than with total integrals.
\end{remark}

\begin{remark}[the sinc kernels are outside this theorem]
The kernels
\begin{equation}
\widehat{w_{T}}(\omega)=\frac{2\sin(T\omega)}{\omega},
\qquad w_{T}=\mathbf{1}_{[-T,T]},
\end{equation}
have total mass $2\pi$ and converge to $2\pi\delta$, but they satisfy none of the hypotheses above:
they change sign infinitely often, and $\int|\widehat{w_{T}}|$ grows like $\log T$. Their convergence
is proved in Chapter~\ref{ch:tespered} by a different route, using the continuity of $\F$ on
$\Sw'$. The existence of this example is the reason Proposition~\ref{prop:general-kernels} is stated
as a sufficient, not a necessary, condition.
\end{remark}

\section{Summary of the chapter}

\begin{center}
\small
\begin{tabular}{@{}p{3.2cm}p{9.6cm}@{}}
\toprule
Statement & Content \\
\midrule
Three modes & pointwise, uniform, distributional; none implies control of integrals
(Examples~\ref{ex:uniform}, \ref{ex:delta-limit}) \\
Licence & dominated convergence, with a \emph{single} integrable bound
(Theorem~\ref{thm:DCT}) \\
Obstruction & no integrable bound exists for $K_{\ep}$
(Proposition~\ref{prop:nodominator}) \\
Sampling & unit mass plus tail concentration implies $\ip{k_{\ep}}{\varphi}\to\varphi(0)$
(Theorem~\ref{thm:sampling-general}) \\
Identity & $\delta*f=f$ (Proposition~\ref{prop:delta-identity}) \\
Necessity & nonnegativity is optional, concentration is not
(Proposition~\ref{prop:general-kernels}) \\
Escape & unit mass alone allows convergence to $0$
(Proposition~\ref{prop:escaping}) \\
\bottomrule
\end{tabular}
\end{center}
